% from ega4-18-fr.tex
\oldpage[IV-4]{123}

de présentation finie, $f$ étant supposé plat au point $x$. Soit $\mathscr J$ l'Idéal quasi-cohérent de type fini de $\sh{O}_S$ qui définit le sous-préschéma fermé $Y$ de $S$. L'hypothèse que $f$ est plat au point $x$ entraîne que l'homomorphisme $\sh{O}_{S,s}\to\sh{O}_{X,x}$ est injectif ($0_{\mathrm I}$, 6.5.1); \emph{a fortiori} l'homomorphisme $\sh{O}_{S,s}\to\sh{O}_{Y,y}$ est injectif, ce qui signifie que $\mathscr J_s=0$, puisque c'est le noyau de l'homomorphisme précédent. Puisque $\mathscr J$ est de type fini, il y a un voisinage ouvert $U$ de $s$ dans $S$ tel que $\mathscr J|U=0$ ($0_{\mathrm I}$, 5.2.2). On peut donc supposer que $h$ est une immersion ouverte, et alors il est clair que $g$ est plat au point $x$.

\begin{proposition}[18.4.10]
\label{IV.18.4.10}
Désignons par $\mathbf P(f,x)$ une propriété vérifiant les conditions suivantes:
\begin{enumerate}
  \item[$1^\circ$] Pour tout morphisme $f:X\to Y$ et tout isomorphisme local $h:Y\to Z$, $\mathbf P(f,x)$ est équivalente à $\mathbf P(h\circ f,x)$ pour $x\in X$.
  \item[$2^\circ$] Pour tout morphisme $f:X\to Y$, tout morphisme étale $g:Y'\to Y$, tout point $x\in X$, si l'on pose $X'=X\times_Y Y'$, $f'=f_{(Y')}:X'\to Y'$ et si $x'\in X'$ est au-dessus de $x$, les propriétés $\mathbf P(f,x)$ et $\mathbf P(f',x')$ sont équivalentes (« invariance par changement de base étale »).
\end{enumerate}
Soient alors $S$ un préschéma, $f:X\to S$ et $h:Y\to S$ deux morphismes, $g:X\to Y$ un $S$-morphisme, $x$ un point de $X$, $y=g(x)$, et supposons $h$ étale au point $y$. Alors les propriétés $\mathbf P(f,x)$ et $\mathbf P(g,x)$ sont équivalentes.
\end{proposition}

Avec les mêmes notations que dans la démonstration de (18.4.9) et en remplaçant $S$ (resp. $Y$) par un voisinage de $s=h(y)$ (resp. de $y$), on peut ici, en vertu de (18.4.8), trouver un morphisme étale $u:S'\to S$ tel que $h'$ soit une immersion ouverte; si $x'\in X'$ est au-dessus de $x$, $\mathbf P(f',x')$ (resp. $\mathbf P(g',x')$) est alors par hypothèse équivalente à $\mathbf P(f,x)$ (resp. $\mathbf P(g,x)$). On peut donc se borner au cas où $h$ est une immersion ouverte. Comme l'hypothèse entraîne que $\mathbf P(g,x)$ est alors équivalente à $\mathbf P(h\circ g,x)$, on en déduit la conclusion.

\begin{env}[18.4.11]
\label{IV.18.4.11}
\emph{Exemples.} --- On peut prendre pour propriété $\mathbf P(f,x)$ l'une quelconque des suivantes, compte tenu de (17.7.4, (ii)):
\begin{enumerate}
  \item[(i)] $f$ est plat au point $x$ (2.2.11, (iv));
  \item[(ii)] $f$ est localement de présentation finie et de coprofondeur $\leq n$ au point $x$ (6.8.1 et 6.7.8);
  \item[(iii)] $f$ est localement de présentation finie et de Cohen-Macaulay au point $x$ (6.8.1 et 6.7.8);
  \item[(iv)] $f$ est localement de présentation finie et possède la propriété $(S_n)$ au point $x$ (6.8.1 et 6.7.8);
  \item[(v)] $f$ est localement de présentation finie et possède la propriété $(R_n)$ au point $x$ (6.8.1 et 6.7.8);
  \item[(vi)] $f$ est localement de présentation finie et normal au point $x$ (6.8.1 et 6.7.8);
  \item[(vii)] $f$ est localement de présentation finie et réduit au point $x$ (6.8.1 et 6.7.8);
  \item[(viii)] $f$ est non ramifié au point $x$ (17.7.4);
  \item[(ix)] $f$ est lisse au point $x$ (17.7.4);
  \item[(x)] $f$ est étale au point $x$ (17.7.4).
\end{enumerate}
\end{env}

